A photoreal 3D copy of a room that a person can walk through needs a model that looks right from anywhere in the room and collision that stops the walker at the real walls. Gaussian splatting provides the first, but plain walls defeat camera poses estimated from images, and the walk collision that viewers build from a splat has not been checked against the real room. We present RoomRecon, which turns an iPhone capture (ARKit poses and LiDAR depth) into a first-person walkable model by reusing maintained tools (nerfstudio with DN-Splatter, splat-transform and the SuperSplat Viewer), and adds the glue they lack and an evaluation of walk collision against a laser scan. On the MuSHRoom \texttt{vr\_room} scene, DN-Splatter reached 26.3\,dB PSNR on held-out frames with a mesh 1.8\,cm from the laser scan (plain splats: 5.9\,cm); our own room reached 22.2\,dB, and both rooms were walkable on a laptop GPU. Collision built from DN-Splatter's mesh matched the scan's walkable area with an IoU of 0.81 and missed 1.7\% of the walls, against 0.58 and 17.8\% for the default collision built from the splat; Polycam's mesh scored 0.79 and 1.8\%. The mesh met three of our four pre-agreed criteria, although a free walk did not reveal the difference. Two other novelty attempts failed their pre-agreed tests and are reported as negative results.
